\chapter*{Как читать это пособие}
\addcontentsline{toc}{chapter}{Как читать это пособие}

Вы уже знаете, что такое цифровой сигнал, умеете читать таблицы истинности
и собирать схемы из элементов И, ИЛИ, НЕ, И-НЕ, ИЛИ-НЕ. Это пособие
продолжает ту же историю: мы научимся собирать \emph{любые} цифровые
схемы не паяльником и не на макетной плате, а \emph{описывая} их текстом, после чего
специальная микросхема (ПЛИС) превращается в ту самую схему, которую
мы описали.

\begin{figure}[H]
\centering
\begin{tikzpicture}[node distance=6mm and 9mm]
  \node[block] (a) {Логические\\элементы};
  \node[block, right=of a] (b) {Триггеры\\и регистры};
  \node[block, right=of b] (c) {Что такое\\ПЛИС};
  \node[hblock, right=of c] (d) {Verilog\\и примеры};
  \node[oblock, below=of d] (e) {Свои\\проекты};
  \node[block, left=of e] (f) {Память\\SDRAM};
  \draw[arr] (a) -- (b); \draw[arr] (b) -- (c); \draw[arr] (c) -- (d);
  \draw[arr] (d) -- (e); \draw[arr] (d) -- (f); \draw[arr] (f) -- (e);
  \node[lbl, above=1mm of a] {вы уже здесь};
\end{tikzpicture}
\caption{Путь, который мы пройдём}
\end{figure}

Пособие построено так:
\begin{enumerate}
  \item \textbf{Главы 1--2} объясняют, что такое ПЛИС, зачем она нужна и что
        находится внутри кристалла.
  \item \textbf{Глава 3} знакомит с платой Sipeed Tang Nano 20K.
  \item \textbf{Глава 4} показывает весь путь от текста до работающей схемы.
  \item \textbf{Глава 5} учит основам языка Verilog.
  \item \textbf{Главы 6--7} содержат шесть законченных проектов и учат
        проверять схему в симуляторе.
  \item \textbf{Глава 8} рассказывает о встроенной памяти SDRAM на 64 Мбит.
  \item \textbf{Глава 9} состоит из заданий для самостоятельной работы.
\end{enumerate}

Все исходные тексты примеров лежат рядом с книгой, в папке
\nolinkurl{personal/fpga/tang_nano_20k/code}. Их можно открыть, скопировать и
сразу загрузить в плату.

\begin{infobox}[Обозначения]
\begin{itemize}
  \item Синие блоки «Справка» содержат полезные подробности.
  \item Оранжевые блоки «Внимание!» предупреждают о типичных ошибках.
  \item Фиолетовые блоки «Главная мысль» подводят итог раздела.
  \item Зелёные блоки «Задание» предлагают что-то сделать самостоятельно.
\end{itemize}
\end{infobox}
